\section{Return-index resolution and its remaining torus complement}
\label{sec:v38-return-resolution}
We apply the preceding estimates to the same actual return law as
Theorem~\ref{thm:LLT}. Write
\begin{equation}\label{eq:v38-return-mass}
 P_{n,m,R}(k,t;J)=
 \nu_R^*\{K_{n,R}=k,\ N_{n,R}=m,\ T_{n,R}-t\in J\}.
\end{equation}
It is zero for $n>m$. The exact disintegration in
Proposition~\ref{prop:v37-exact-return-disintegration} applies to a
single $n$ and to every block of consecutive $n$; no limiting step is
needed for either statement.

\begin{theorem}[Concentration and fine windows for actual returns]
\label{thm:v38-fine-return-windows}
For all $m\ge n\ge1$, all $k,t,R$ and bounded intervals $J$,
\begin{equation}\label{eq:v38-actual-singleton-upper}
 P_{n,m,R}(k,t;J)\le C(1+|J|)m^{-2}.
\end{equation}
For a fixed $J$ of positive length, a block
$\mathcal B=\{j,\ldots,j+H-1\}$ with center $\ell=j+(H-1)/2$, and
$Z=(z,y)$ as in \eqref{eq:v38-central-coordinates}, one has
\begin{equation}\label{eq:v38-fine-return-window}
 \frac{m^2}{H}\sum_{n\in\mathcal B\cap[1,m]}P_{n,m,R}(k,t;J)
       -c|J|\mathcal G_{m,H,R}(z,y)\longrightarrow0,
\end{equation}
uniformly for $R$, $|Z|\le M$, and $H\ge h_m$, for any $h_m\to\infty$.
If $h_m\le H\le u_m\sqrt m$ and $u_m\to0$, this becomes
\begin{equation}\label{eq:v38-subdiffusive-return}
 \frac{m^2}{H}\sum_{n\in\mathcal B}P_{n,m,R}(k,t;J)
       =c|J|g_{\Omega_R}(z,y)+o(1).
\end{equation}
The block lies in $[1,m]$ for all sufficiently large $m$ in this last
central regime. Its probability is bounded below by $c_{M,J}Hm^{-2}$.
The same lower bound holds for $H\le C_0\sqrt m$ with a constant
also depending on $C_0$.
\end{theorem}
\begin{proof}
Use $a_R=d_R=\eta_R$ in
Theorems~\ref{thm:v38-singleton-concentration} and
\ref{thm:v38-occupation-mesoscopic}. These endpoints are admissible
by Lemma~\ref{lem:v37-section-multiplier}. Their amplitude is $c^2$.
The exact disintegration divides the collision integral by $c$, so
the amplitude is $c$, proving every displayed formula. Since
$0<c<1$, a block centered at $mc+O(\sqrt m)$ and of width
$O(\sqrt m)$ is contained in $[1,m]$ eventually. On the compact set
$|Z|\le M$ and $|s|/\sqrt m\le C_0/2$, uniform ellipticity gives a
positive lower bound for $g_{\Omega_R}(z,y+s/\sqrt m)$. The error
in \eqref{eq:v38-fine-return-window} tends uniformly to zero.
\end{proof}

\begin{corollary}[The original return covariance normalization]
\label{cor:v38-return-normalization}
In the subdiffusive regime of \eqref{eq:v38-subdiffusive-return}, put
\[
 V_{\ell,R}(k,m,t)
  =\ell^{-1/2}(k_1,k_2,m-\ell/c,t-\ell\bar\tau_R/c).
\]
Then, uniformly on the same central compact sets,
\begin{equation}\label{eq:v38-original-normalization}
 \frac{\ell^2}{H}\sum_{n\in\mathcal B}P_{n,m,R}(k,t;J)
       =|J|g_{D_R}\bigl(V_{\ell,R}(k,m,t)\bigr)+o(1).
\end{equation}
The center $\ell$ may be a half-integer; every summand still has an
integer return index.
\end{corollary}
\begin{proof}
Equation~\eqref{eq:v37-return-gaussian-agreement} gives
$c g_{\Omega_R}(Z)=c^{-2}g_{D_R}(c^{-1/2}L_R^{-1}Z)$.
Here $L_R^{-1}Z=m^{-1/2}(k_1,k_2,m-\ell/c,t-\ell\bar\tau_R/c)$
exactly, and $\ell/m=c+O(m^{-1/2})$. Multiply
\eqref{eq:v38-subdiffusive-return} by $\ell^2/m^2$ and use uniform
continuity of $g_{D_R}$ on compact sets. No independent covariance
or clock is introduced.
\end{proof}

\begin{corollary}[Regular endpoint amplitudes in a fine return window]
\label{cor:v38-selected-return-window}
Let $a_R,d_R$ be bounded nonnegative functions such that
$\eta_Ra_R,\eta_Rd_R$ form multiplier-bounded endpoints.
Insert $a_R(x)d_R((F_R^*)^nx)$ in each original return probability.
Then the right side of \eqref{eq:v38-fine-return-window} is multiplied
by $\nu_R^*(a_R)\nu_R^*(d_R)$. A uniform positive lower bound for the
selected probability requires this product to be uniformly positive.
The concentration estimate holds with a changed endpoint constant.
\end{corollary}
\begin{proof}
On $N_{n,R}=m$ one has $(F_R^*)^nx=T_R^mx$. Apply the two collision
estimates to $\eta_Ra_R,\eta_Rd_R$. Their product of collision means,
divided by $c$, equals $c\nu_R^*(a_R)\nu_R^*(d_R)$.
\end{proof}

\subsection{Exact single-index inversion with a localized unresolved term}
The upper bound and the window asymptotic do not alone identify each
single return mass. The following identity locates the remaining
condition directly on the bounded collision operator, rather than on
an unspecified induced spectral realization.

Let $q$ be an integrable roof test with integrable Fourier transform
supported in $[-B,B]$. Define the actual single-return functional
\begin{equation}\label{eq:v38-band-return-functional}
 P^q_{n,m,R}(k,t)=\int_{Y_R^*}
 \one_{\{K_{n,R}=k,N_{n,R}=m\}}q(T_{n,R}-t)\,d\nu_R^*.
\end{equation}
In the following formulas $\Phi^{\eta,\eta}_{m,R}$ has the two actual
section endpoints. Choose a real smooth function $\chi_B$ on the
occupation torus, equal to one near zero and supported in
$(-\sigma_B,\sigma_B)$. It is chosen once after $B$.
Set
\begin{align}\label{eq:v38-fixed-torus-remainder}
 \mathcal R^q_{n,m,R}(k,t)=\frac1{c(2\pi)^4}
 \int_{[-\pi,\pi]^2}\int_{-B}^{B}\int_{-\pi}^{\pi}
 &e^{-i[u\cdot k+b(t-m\bar\tau_R)+v(n-mc)]}
 \widehat q(-b)(1-\chi_B(v))\notag\\[-2mm]
 &\hspace{-10mm}\cdot\Phi^{\eta,\eta}_{m,R}(u,b,v)\,dv\,db\,du.
\end{align}
Every operator in this expression is a finite power of bounded
one-step multipliers. No uniform bound on its long powers over the
whole occupation torus is presupposed.

\begin{proposition}[Exact-index Gaussian term and the fixed-torus remainder]
\label{prop:v38-exact-index-remainder}
With $Z=(k/\sqrt m,(t-m\bar\tau_R)/\sqrt m,(n-mc)/\sqrt m)$,
\begin{equation}\label{eq:v38-exact-index-remainder}
 m^2P^q_{n,m,R}(k,t)
 =c\left(\int q\right)g_{\Omega_R}(Z)
          +m^2\mathcal R^q_{n,m,R}(k,t)+o(1),
\end{equation}
uniformly in $R$ and central compact sets, for each fixed $q,B,\chi_B$.
Thus for this roof test the single-index local limit is equivalent to
$m^2\mathcal R^q_{n,m,R}(k,t)\to0$ in the same uniform sense.
\end{proposition}
\begin{proof}
Apply the exact return disintegration, torus orthogonality for
$K_m^{\rm c}$ and the integer $A_m$, and real Fourier inversion for
$q(S_m-t)$. This gives the full integral in
\eqref{eq:v38-fixed-torus-remainder} with $1-\chi_B$ replaced by one.
Centering the occupation only moves a scalar phase between the inverse
factor and $\Phi$; their product remains $2\pi$-periodic in $v$.
The identity is therefore valid even when $mc$ is not an integer.

Split the integral with $1=\chi_B+(1-\chi_B)$. On the first piece,
Proposition~\ref{prop:v38-strip-integral} and the joint expansion
allow the four-dimensional rescaling $w=\sqrt m(u,b,v)$.
The factors $\chi_B(w_4/\sqrt m)$ and
$\widehat q(-w_3/\sqrt m)$ tend to one and $\int q$, respectively.
The endpoint amplitude is $c^2$, divided by the source normalization
$c$. Gaussian domination proves the first term and the uniform $o(1)$.
The second piece is exactly the stated remainder. Subtraction proves
the equivalence, with no estimate of the modulus of its integrand.
\end{proof}

\begin{corollary}[A sufficient occupation-torus condition for a sharp roof interval]
\label{cor:v38-single-index-criterion}
Fix a roof interval $J$. Suppose that, for every fixed $B\ge1$ and
its two interval envelopes $q_{J,B}^{\pm}$,
\[
 \sup_{R,\,|Z|\le M}m^2
       |\mathcal R^{q_{J,B}^{\pm}}_{n,m,R}(k,t)|\longrightarrow0.
\]
Then $m^2P_{n,m,R}(k,t;J)=c|J|g_{\Omega_R}(Z)+o(1)$ uniformly on
those central compact sets. This is a sufficient condition, not an
additional proved cancellation assertion.
\end{corollary}
\begin{proof}
Multiply the two roof-envelope inequalities by the unchanged positive
single-index source measure. For each fixed $B$ the proposition and
the assumed cancellation give lower and upper limiting Gaussian
masses $c(|J|-2\pi/B)$ and $c(|J|+2\pi/B)$. Take $m\to\infty$
first and $B\to\infty$ afterwards.
\end{proof}

\paragraph{Proof dependencies and relation to the raw endpoint.}
The new unconditional results use the existing exact occupation
multiplier, joint covariance and small-occupation resolvent estimate,
together with the one-dimensional interval envelopes. They add no
continuum partition estimate and do not assume a gap on a full
occupation torus. The fixed-strip argument is a local-limit and
concentration mechanism; its present application concerns the actual
five-rectangle section and its true return law. The classical Fourier
envelopes are not claimed as new, and a general priority claim for
mesoscopic local limits is not made.

Compared with Theorem~\ref{thm:v37-return-window}, the window may now
contain $H_m$ consecutive return counts with $H_m\to\infty$ arbitrarily
slowly, rather than a fixed positive multiple of $\sqrt m$.
Equation~\eqref{eq:v38-actual-singleton-upper} also applies to each
single return index at its natural four-dimensional upper scale.
Neither result sets $H_m=1$ in the window limit: the explicit error
$C_J/(\sigma H_m)$ in \eqref{eq:v38-two-band-budget} would persist.
The cancellation in \eqref{eq:v38-fixed-torus-remainder} is the
separate exact-index requirement. Pointwise roof inversion additionally
requires the common raw correction and full roof complement of the
original four-coordinate theorem; fixed roof intervals do not supply
it. All original raw-edge and return-inversion statements are retained.
